\documentclass[11pt,a4paper]{article}

\usepackage[utf8]{inputenc}
\usepackage[T1]{fontenc}
\usepackage[french]{babel}
\usepackage{lmodern}
\usepackage[margin=2cm]{geometry}
\usepackage{xcolor}
\usepackage{tikz}
\usepackage{booktabs}
\usepackage{tabularx}
\usepackage{array}
\usepackage{siunitx}
\usepackage{amsmath}
\usepackage{amssymb}
\usepackage{pifont}
\usepackage{enumitem}
\usepackage[colorlinks=true, linkcolor=blue!50!black, urlcolor=blue!60!black]{hyperref}

\usetikzlibrary{arrows.meta, positioning, shapes.geometric, calc, fit, backgrounds, decorations.pathreplacing}

% Couleurs fonctionnelles
\definecolor{chw}{HTML}{B71C1C}      % hardware (rouge foncé)
\definecolor{cdsp}{HTML}{1565C0}     % DSP (bleu)
\definecolor{ca53}{HTML}{2E7D32}     % A53 userspace (vert)
\definecolor{cusb}{HTML}{6A1B9A}     % USB (violet)
\definecolor{cnpu}{HTML}{E65100}     % NPU (orange)
\definecolor{cfx}{HTML}{00838F}      % FX (cyan)
\definecolor{cbg}{HTML}{F5F5F5}

% Styles de blocs
\tikzset{
  hw/.style={draw, fill=chw!15, rounded corners=2pt, font=\sffamily\footnotesize, align=center, minimum height=7mm, inner sep=3pt},
  dsp/.style={draw, fill=cdsp!15, rounded corners=2pt, font=\sffamily\footnotesize, align=center, minimum height=7mm, inner sep=3pt},
  a53/.style={draw, fill=ca53!15, rounded corners=2pt, font=\sffamily\footnotesize, align=center, minimum height=7mm, inner sep=3pt},
  usb/.style={draw, fill=cusb!15, rounded corners=2pt, font=\sffamily\footnotesize, align=center, minimum height=7mm, inner sep=3pt},
  npu/.style={draw, fill=cnpu!15, rounded corners=2pt, font=\sffamily\footnotesize, align=center, minimum height=7mm, inner sep=3pt},
  fx/.style={draw, fill=cfx!15, rounded corners=2pt, font=\sffamily\footnotesize, align=center, minimum height=7mm, inner sep=3pt},
  zone/.style={draw, dashed, rounded corners=6pt, inner sep=10pt},
  arr/.style={-{Stealth[length=2mm]}, thick},
  arrmulti/.style={-{Stealth[length=2mm]}, very thick},
  pcm/.style={draw, fill=yellow!20, font=\sffamily\scriptsize, minimum height=5mm, rounded corners=1pt}
}

\title{\textbf{Plateforme Mastering Live Debix}\\[4pt]
  \large Architectures DSP --- comparatif de 3 propositions}
\author{Document de travail}
\date{2026-04-20}

\begin{document}
\maketitle

\tableofcontents
\vspace{1em}

\hrule
\vspace{0.5em}
\noindent\textit{Ce document compare \textbf{3 architectures DSP} pour la plateforme mastering live. Objectif~: te permettre de trancher un choix structurant avant tout code. Chaque architecture a un schéma, un inventaire de PCMs, un budget latence, un budget CPU, ses avantages et ses contraintes.}
\vspace{0.5em}
\hrule

% =======================================================================
\section{Contexte et rappel}

\subsection{Matériel}
\begin{itemize}[leftmargin=*, itemsep=2pt]
  \item \textbf{SoC}~: NXP i.MX8MP --- 4× Cortex-A53 @ 1{,}6~GHz + Cortex-M7 + DSP HiFi4 @ 800~MHz + NPU Vivante VIP8000 @ 2{,}3~TOPS
  \item \textbf{Front-end audio}~: 4× codec TAC5212 daisy-chaînés sur SAI7, TDM 8~canaux $\times$ 32~bits @ 48~kHz
  \item \textbf{USB gadget}~: 2 contrôleurs USB $\rightarrow$ UAC2 8in/8out (ASIO DAW) + UAC2 2in/2out (iOS)
\end{itemize}

\subsection{Objectif fonctionnel}
Faire un \textbf{mastering live "lite"} piloté par le NPU. Le NPU analyse le bus master stéréo en continu et ajuste les paramètres (EQ, multiband compressor, exciter) à $\sim$10~Hz d'update.

\subsection{Contraintes}
\begin{center}
\begin{tabularx}{0.92\textwidth}{lX}
\toprule
\textbf{Contrainte} & \textbf{Valeur} \\
\midrule
Sample rate & 48~kHz fixe (Phase 1) \\
Latence chemin principal (USB\,$\rightarrow$\,DSP\,$\rightarrow$\,USB) & \textbf{$\leq$ 10~ms} \\
Latence effets send (reverb, chorus\ldots) & non critique (20--50~ms OK) \\
DSP period SOF & 2~ms (pipeline scheduling) \\
Effets send & \textbf{en userspace A53} (JACK + LV2), pas dans le DSP \\
NPU & modèle TFLite à entraîner \emph{from scratch} \\
Règle & minimiser le contenu du DSP \\
\bottomrule
\end{tabularx}
\end{center}

\subsection{Les 4 axes qui distinguent les architectures}
\label{sec:axes}

Chaque architecture répond différemment à ces 4 questions~:

\begin{enumerate}[leftmargin=*, itemsep=3pt]
  \item \textbf{Où vit la matrice de routage ?} (DSP ou A53 userspace)
  \item \textbf{Où vit l'EQ fine des entrées ?} (Codec TAC5212 / DSP / A53)
  \item \textbf{Où vit la chaîne master ?} (DSP ou A53)
  \item \textbf{Comment les audio USB transitent-elles entre le gadget et le DSP ?} (passage direct par A53 kernel / bridge \texttt{alsaloop} / JACK)
\end{enumerate}

\begin{center}
\begin{tabular}{lccc}
\toprule
\textbf{Axe} & \textbf{Archi A (Max DSP)} & \textbf{Archi B (Min DSP)} & \textbf{Archi C (Équilibrée)} \\
\midrule
Matrice de routage & DSP & A53 (JACK) & DSP \\
EQ fine entrées & DSP (\texttt{eq\_iir}) & A53 (LV2) & Codec + DSP léger \\
Chaîne master & DSP & A53 & DSP \\
USB$\leftrightarrow$DSP bridge & \texttt{alsaloop} & JACK natif & \texttt{alsaloop} \\
\bottomrule
\end{tabular}
\end{center}

\clearpage
% =======================================================================
\section{Architecture A --- Maximum DSP}
\label{sec:archiA}

\textbf{Philosophie}~: tout ce qui peut être fait en DSP y est fait. L'A53 sert uniquement de plombier ALSA vers USB, héberge les 4 effets send (obligatoire --- SOF n'a pas reverb/chorus/flanger), et fait tourner le NPU et l'UI.

\subsection{Schéma}

\begin{center}
\begin{tikzpicture}[node distance=4mm and 8mm]

% --- Hardware inputs
\node[hw]                 (mics)  {8 mics};
\node[hw, below=of mics]  (tac)   {4× TAC5212\\(AGC + HPF\\+ 3 biquads/ch)};

% --- DSP block
\node[dsp, right=18mm of tac, minimum width=70mm, minimum height=60mm, align=center] (dsp) {};
\node[anchor=north] at (dsp.north) {\textbf{\textcolor{cdsp}{DSP HiFi4 (SOF)}}};

\node[dsp, fill=cdsp!25] (eq8)     at ($(dsp.west)+(1.3,1.8)$) {8× eq\_iir\\(fine EQ)};
\node[dsp, fill=cdsp!25, right=5mm of eq8] (matrix) {Matrice\\18$\times$18\\mixin\_mixout};
\node[dsp, fill=cdsp!25, below=5mm of matrix] (mbdrc)  {multiband\_drc};
\node[dsp, fill=cdsp!25, below=3mm of mbdrc]  (vol)    {volume};
\node[dsp, fill=cdsp!25, below=3mm of vol]    (lim)    {drc (limiter)};
\node[dsp, fill=cdsp!25, left=5mm of mbdrc]   (tap)    {NPU tap};

% --- A53 block
\node[a53, right=22mm of dsp, minimum width=55mm, minimum height=60mm] (a53) {};
\node[anchor=north] at (a53.north) {\textbf{\textcolor{ca53}{A53 userspace}}};

\node[fx, fill=cfx!25] (jack) at ($(a53.west)+(1.2,1.6)$) {JACK\\+ 4 LV2\\(Reverb\\Chorus\\Flanger\\Delay)};
\node[usb, fill=cusb!25, below=3mm of jack] (usb1) {alsaloop\\USB1 (8in/8out)};
\node[usb, fill=cusb!25, below=3mm of usb1] (usb2) {alsaloop\\USB2 (2in/2out)};
\node[npu, fill=cnpu!25, below=3mm of usb2] (npu) {NPU service};

% --- USB outside
\node[usb, right=8mm of usb1] (daw)  {DAW\\ASIO};
\node[usb, right=8mm of usb2] (ios)  {iOS};

% --- Speakers
\node[hw, left=12mm of lim] (hp) {HP/monitor};

% Arrows: hardware to DSP
\draw[arrmulti] (mics) -- (tac);
\draw[arrmulti] (tac.east) -- node[above, font=\scriptsize, align=center]{SAI7\\8ch TDM} (eq8.west);

% Internal DSP
\draw[arr] (eq8) -- (matrix);
\draw[arr] (matrix.south) -- (mbdrc.north);
\draw[arr] (mbdrc) -- (vol);
\draw[arr] (vol) -- (lim);
\draw[arr] (mbdrc) -- (tap);

% DSP to HP (via matrix 8ch direct)
\draw[arrmulti, dashed] (matrix.west) to[bend right=40] node[left, font=\scriptsize]{8ch out} (hp);
\draw[arr] (lim.west) to[bend left=30] (hp);

% DSP <-> A53
\draw[arrmulti, <->] (matrix.east)+(0,2mm) -- node[above, font=\scriptsize, align=center]{FX 4×2ch\\send/return} (jack.west);
\draw[arrmulti, <->] (matrix.east)+(0,-2mm) -- node[below, font=\scriptsize, align=center, pos=0.3]{USB 8+2\\send/return} (usb1.west);
\draw[arr] (tap.west) to[out=180, in=200] ($(dsp.south west)+(5mm,0)$) to node[below=2mm, font=\scriptsize]{analysis} (npu.west);

% A53 to USB outside
\draw[arrmulti, <->] (usb1.east) -- (daw.west);
\draw[arrmulti, <->] (usb2.east) -- (ios.west);

% NPU control back
\draw[arr, dashed, red!70!black] (npu.north west) to[bend left=60] node[right=2mm, font=\scriptsize, text=red!70!black]{ALSA ctrl} (matrix.south east);

\end{tikzpicture}
\end{center}

\subsection{Inventaire des PCMs (côté A53 = côté kernel ALSA)}
\begin{center}
\begin{tabularx}{\textwidth}{lllX}
\toprule
\textbf{PCM} & \textbf{Type} & \textbf{Ch} & \textbf{Rôle} \\
\midrule
\texttt{SAI\_Dry} & capture & 8 & Mics bruts post-TAC5212 (bypass matrice) \\
\texttt{Master} & capture & 2 & Sortie master post-chaîne (NPU + bridge USB) \\
\texttt{USB1\_Send} & capture & 8 & Sortie matrice vers DAW \\
\texttt{USB2\_Send} & capture & 2 & Sortie matrice vers iOS \\
\texttt{FX1..4\_Send} & capture $\times$4 & 2 & Sorties vers JACK \\
\texttt{SAI\_Monitor} & playback & 8 & Vers TAC5212 DAC (HP/monitor) \\
\texttt{USB1\_Ret} & playback & 8 & Retour depuis DAW \\
\texttt{USB2\_Ret} & playback & 2 & Retour depuis iOS \\
\texttt{FX1..4\_Ret} & playback $\times$4 & 2 & Retours depuis JACK \\
\bottomrule
\end{tabularx}
\end{center}
\textbf{Total~: 12 PCMs} (4 captures fixes + 4 captures FX + 3 playback fixes + 4 playback FX = 15 pipelines SOF).

\subsection{Budget latence (chemin principal USB$\rightarrow$DSP$\rightarrow$USB)}
\begin{center}
\begin{tabular}{lr}
\toprule
\textbf{Étape} & \textbf{Latence} \\
\midrule
USB gadget UAC2 in (isochronous, 48 kHz) & 1{,}0~ms \\
\texttt{alsaloop} gadget$\rightarrow$DSP PCM (période 2~ms, 2 périodes) & 4{,}0~ms \\
DSP pipeline 2~ms (1 cycle in + 1 cycle out + matrice) & 4{,}0~ms \\
\texttt{alsaloop} DSP PCM$\rightarrow$gadget & 4{,}0~ms \\
USB gadget UAC2 out & 1{,}0~ms \\
\midrule
\textbf{Total round-trip} & \textbf{14{,}0~ms} \\
\bottomrule
\end{tabular}
\end{center}
\textcolor{red!70!black}{\textbf{$>$ 10~ms~!}} Dépasse la contrainte. Nécessite DSP period 1~ms (au lieu de 2~ms) et \texttt{alsaloop} très agressif, ce qui stresse le HiFi4.

\subsection{Budget CPU DSP (estimation)}
\begin{center}
\begin{tabular}{lr}
\toprule
\textbf{Composant} & \textbf{\% HiFi4} \\
\midrule
8 $\times$ eq\_iir (3 biquads/ch) & $\sim$5~\% \\
Matrice 18$\times$18 mixin\_mixout & $\sim$8~\% \\
multiband\_drc stéréo 3 bandes & $\sim$12~\% \\
volume + drc limiter & $\sim$2~\% \\
12 pipelines + IPC overhead & $\sim$8~\% \\
\midrule
\textbf{Total} & $\sim$35~\% \\
\bottomrule
\end{tabular}
\end{center}
CPU confortable, mais topology \textbf{très complexe} (15 pipelines, 4 effets send à câbler en PCM loopback, chaque FX = 2 PCMs + 2 pipelines).

\subsection{Avantages et inconvénients}
\begin{center}
\begin{tabularx}{\textwidth}{p{0.48\textwidth} X}
\toprule
\textbf{\textcolor{ca53}{Avantages}} & \textbf{\textcolor{chw}{Inconvénients}} \\
\midrule
Routage 100~\% déterministe, non dépendant de la charge Linux & 15 pipelines SOF à maintenir~: topology $\sim$500 lignes m4 \\
Même latence pour tous les chemins & Latence round-trip $>$ 10~ms (contrainte violée) \\
Paramètres matrix visibles via \texttt{amixer} & Chaque ajout (5\textsuperscript{e} FX, nouvel input) $=$ refonte topology + rebuild firmware \\
CPU DSP sous-exploité & Tous les bugs sont dans le DSP, debug plus difficile \\
& Matrice 18$\times$18 $=$ 324 cellules$\times$contrôles ALSA $=$ bruit dans l'UI \\
\bottomrule
\end{tabularx}
\end{center}

\subsection{Verdict Archi A}
\textbf{Non retenue} si on veut tenir 10~ms. Trop complexe pour le ratio coût/bénéfice.

% =======================================================================
\clearpage
\section{Architecture B --- Minimum DSP ("Thin DSP")}
\label{sec:archiB}

\textbf{Philosophie}~: le DSP est réduit au strict minimum --- juste un adaptateur codec + un tap pour le NPU. Tout le routage, l'EQ, le master chain et les effets send vivent dans un \textbf{graphe JACK} sur l'A53. Le codec TAC5212 fait le pré-conditionnement (AGC, biquads, HPF) pour soulager le DSP.

\subsection{Schéma}

\begin{center}
\begin{tikzpicture}[node distance=4mm and 8mm]

% Hardware
\node[hw]                 (mics)  {8 mics};
\node[hw, below=of mics]  (tac)   {4× TAC5212\\(AGC + HPF\\+ 3 biquads/ch\\+ phase/gain cal)};

% DSP thin
\node[dsp, right=15mm of tac, minimum width=40mm, minimum height=35mm] (dsp) {};
\node[anchor=north] at (dsp.north) {\textbf{\textcolor{cdsp}{DSP HiFi4}}};
\node[dsp, fill=cdsp!25] (passin)  at ($(dsp.west)+(1.2,0.5)$) {SAI$\leftrightarrow$host\\pass-through 8ch};
\node[dsp, fill=cdsp!25, below=2mm of passin] (tap) {Master tap\\2ch};

% A53 big
\node[a53, right=18mm of dsp, minimum width=75mm, minimum height=75mm] (a53) {};
\node[anchor=north] at (a53.north) {\textbf{\textcolor{ca53}{A53 userspace (JACK graph)}}};

\node[fx, fill=ca53!25] (eq) at ($(a53.west)+(1.2,2.8)$) {8× eq LV2\\(fine EQ)};
\node[fx, fill=ca53!25, right=4mm of eq] (matx) {Routing matrix\\JACK internal\\(N$\times$M)};
\node[fx, fill=ca53!25, below=3mm of matx] (master) {Master chain\\mbdrc $\rightarrow$ vol\\$\rightarrow$ limiter\\(LV2)};
\node[fx, fill=ca53!25, left=4mm of master] (fxall) {4 LV2 FX\\reverb\\chorus\\flanger\\delay};

\node[usb, fill=cusb!25, below=3mm of fxall] (usb1) {UAC2 gadget\\USB1 (8+8)};
\node[usb, fill=cusb!25, right=4mm of usb1] (usb2) {UAC2 gadget\\USB2 (2+2)};
\node[npu, fill=cnpu!25, right=4mm of usb2] (npu) {NPU};

% External USB clients
\node[usb, right=8mm of usb2] (daw)  {DAW\\ASIO};
\node[usb, above right=-2mm and 8mm of daw] (ios)  {iOS};

\node[hw, left=10mm of master] (hp) {HP/monitor};

% Arrows
\draw[arrmulti] (mics) -- (tac);
\draw[arrmulti] (tac.east) -- node[above, font=\scriptsize]{SAI7 TDM 8ch} (passin.west);
\draw[arrmulti, <->] (passin.east) -- node[above, font=\scriptsize]{ALSA 8+8} (eq.west);
\draw[arr] (eq) -- (matx);
\draw[arr] (matx.south) -- (master.north);
\draw[arr, <->] (matx.west) to[bend right=15] (fxall.north);
\draw[arr] (master) -- node[above, font=\scriptsize]{loopback} (tap.east);
\draw[arr] (master.west) -- (hp);
\draw[arrmulti, <->] (matx.south east) to[bend left=25] (usb1.north);
\draw[arrmulti, <->] (matx.east) to[bend left=10] (usb2.north);
\draw[arrmulti, <->] (usb1) -- (daw);
\draw[arrmulti, <->] (usb2) -- (ios);
\draw[arr, dashed, red!70!black] (tap.south east) to[bend right=20] (npu.west);
\draw[arr, dashed, red!70!black] (npu.north) to[bend right=10] node[right, font=\scriptsize, text=red!70!black]{ALSA ctrl} (master.south east);

\end{tikzpicture}
\end{center}

\subsection{Inventaire des PCMs (très court)}
\begin{center}
\begin{tabular}{lllp{0.42\textwidth}}
\toprule
\textbf{PCM} & \textbf{Type} & \textbf{Ch} & \textbf{Rôle} \\
\midrule
\texttt{SAI\_Capture} & capture & 8 & Mics TAC5212 vers JACK \\
\texttt{SAI\_Playback} & playback & 8 & JACK vers TAC5212 DAC \\
\texttt{Master\_Tap} & capture & 2 & Master post-chain (NPU) \\
\bottomrule
\end{tabular}
\end{center}
\textbf{Total~: 3 PCMs côté DSP}. Tout le reste vit dans JACK.

\subsection{Budget latence}
\begin{center}
\begin{tabular}{lr}
\toprule
\textbf{Étape} & \textbf{Latence} \\
\midrule
USB gadget in & 1{,}0~ms \\
JACK buffer (period 128 frames @ 48k $=$ 2{,}7~ms, $\times$2) & 5{,}3~ms \\
LV2 plugins (EQ + matrix + master chain) & $\sim$2~ms \\
USB gadget out & 1{,}0~ms \\
\midrule
\textbf{Round-trip sans DSP} & \textbf{9{,}3~ms} \\
\midrule
\textit{Chemin avec DSP (mic$\rightarrow$DAC)~:} & \\
SAI in (1 DSP period 2ms, contrôle DMA) & 2{,}0~ms \\
DSP$\rightarrow$host (ALSA 1 période de 2ms) & 2{,}0~ms \\
JACK graph & 5{,}3~ms \\
host$\rightarrow$DSP (ALSA 2ms) & 2{,}0~ms \\
SAI out & 2{,}0~ms \\
\textbf{Total mic$\rightarrow$HP} & \textbf{13{,}3~ms} \\
\bottomrule
\end{tabular}
\end{center}
\textcolor{red!70!black}{\textbf{13 ms pour mic$\rightarrow$HP}}. Mais le chemin "principal" USB$\rightarrow$USB reste à $\sim$9~ms si l'audio ne traverse pas le DSP (tout JACK). En pratique on veut que les mics traversent quand même le DSP pour le pré-conditionnement et le tap NPU, donc \textbf{on sort du budget 10~ms pour le mic monitoring live}.

\subsection{Budget CPU}
Charge DSP $\approx$ 3~\% (juste DMA et pipelines pass-through).
Charge A53 $\approx$ 35--45~\% (JACK + 4 FX + EQ + matrix + master), soit \textbf{quasi 2 cœurs sur 4} à pleine charge.

\subsection{Avantages et inconvénients}
\begin{center}
\begin{tabularx}{\textwidth}{p{0.48\textwidth} X}
\toprule
\textbf{\textcolor{ca53}{Avantages}} & \textbf{\textcolor{chw}{Inconvénients}} \\
\midrule
Topology SOF \textbf{triviale} ($<$ 50 lignes m4) & Latence mic$\rightarrow$HP $>$ 10~ms (contrainte violée pour monitoring live) \\
Écosystème LV2 mature (reverb, mastering plugins, éditeurs) & Charge A53 élevée --- compétition avec Flutter UI et NPU \\
Plugins hot-swappables à chaud sans rebuild firmware & Dépendant de la stabilité du scheduler Linux (xruns possibles sous charge) \\
Graphe JACK visible avec patchbay (catia/qjackctl) & 2 cœurs A53 consommés en pic $=$ moins de marge pour NPU \\
Debug plus simple (outils Linux standards) & Bridge ALSA DSP$\leftrightarrow$JACK ajoute 2--4~ms \\
\bottomrule
\end{tabularx}
\end{center}

\subsection{Verdict Archi B}
\textbf{Adaptée si} la latence live n'est pas critique ou si on accepte un monitoring direct via le codec (bypass analogique TAC5212). \textbf{Rejetée} ici car contrainte 10~ms non tenue sur le chemin principal.

% =======================================================================
\clearpage
\section{Architecture C --- Équilibrée (recommandée)}
\label{sec:archiC}

\textbf{Philosophie}~: le DSP fait ce qui \textbf{DOIT} être fait à basse latence (matrice de routage + chaîne master stéréo + tap NPU) et rien d'autre. Le codec TAC5212 fait le pré-conditionnement par canal. Les effets send (non critiques) vivent dans JACK/LV2 sur A53. La répartition joue sur les forces de chaque étage.

\subsection{Schéma}

\begin{center}
\begin{tikzpicture}[node distance=4mm and 8mm]

% Hardware
\node[hw]                 (mics)  {8 mics};
\node[hw, below=of mics]  (tac)   {4× TAC5212\\(AGC + HPF\\+ 3 biquads/ch\\+ phase/gain cal)};

% DSP
\node[dsp, right=14mm of tac, minimum width=55mm, minimum height=62mm] (dsp) {};
\node[anchor=north] at (dsp.north) {\textbf{\textcolor{cdsp}{DSP HiFi4 (SOF)}}};

\node[dsp, fill=cdsp!25] (eq) at ($(dsp.west)+(1.1,2.0)$) {8× eq\_iir\\(touche fine)};
\node[dsp, fill=cdsp!25, below=2.5mm of eq] (matrix) {Matrice\\mixin\_mixout};
\node[dsp, fill=cdsp!25, below=2.5mm of matrix] (mbdrc) {multiband\_drc\\(L/R, 3 bandes)};
\node[dsp, fill=cdsp!25, below=2mm of mbdrc] (voldrc) {volume $\rightarrow$ drc\\(limiter)};
\node[dsp, fill=cdsp!25, right=4mm of mbdrc] (tap) {NPU\\tap};

% A53
\node[a53, right=18mm of dsp, minimum width=55mm, minimum height=62mm] (a53) {};
\node[anchor=north] at (a53.north) {\textbf{\textcolor{ca53}{A53 userspace}}};

\node[fx, fill=cfx!25] (jack) at ($(a53.west)+(1.2,2.0)$) {JACK + 4 LV2\\Reverb\\Chorus\\Flanger\\Delay};
\node[usb, fill=cusb!25, below=3mm of jack] (usb1) {UAC2 gadget\\USB1 (8in/8out)\\alsaloop bridge};
\node[usb, fill=cusb!25, below=3mm of usb1] (usb2) {UAC2 gadget\\USB2 (2in/2out)};
\node[npu, fill=cnpu!25, below=3mm of usb2] (npu) {NPU + UI\\(Flutter)};

\node[usb, right=8mm of usb1] (daw) {DAW\\ASIO};
\node[usb, right=8mm of usb2] (ios) {iOS};

\node[hw, left=10mm of voldrc] (hp) {HP/monitor};

% Arrows in DSP chain
\draw[arrmulti] (mics) -- (tac);
\draw[arrmulti] (tac.east) -- node[above, font=\scriptsize, align=center]{SAI7\\TDM 8ch} (eq.west);
\draw[arr] (eq) -- (matrix);
\draw[arr] (matrix) -- (mbdrc);
\draw[arr] (mbdrc) -- (voldrc);
\draw[arr] (mbdrc.east) -- (tap.west);

% DSP output to speakers
\draw[arrmulti, dashed] (matrix.west) to[bend right=45] node[left, font=\scriptsize]{8ch direct} (hp);
\draw[arr] (voldrc.west) -- (hp);

% A53 traffic
\draw[arrmulti, <->] (matrix.east) -- node[above, font=\scriptsize, align=center]{FX 4×2ch\\ALSA} (jack.west);
\draw[arrmulti, <->] (matrix.east)+(0,-3mm) -- node[below, font=\scriptsize, align=center]{USB 8+2\\ALSA} (usb1.west);
\draw[arrmulti, <->] (usb1.east) -- (daw.west);
\draw[arrmulti, <->] (usb2.east) -- (ios.west);
\draw[arr, dashed, red!70!black] (tap.south) to[bend left=10] (npu.north west);
\draw[arr, dashed, red!70!black] (npu.north) to[bend right=15] node[right=2mm, font=\scriptsize, text=red!70!black]{ctrls} (matrix.south east);

\end{tikzpicture}
\end{center}

\subsection{Inventaire PCMs}
\begin{center}
\begin{tabularx}{\textwidth}{llX}
\toprule
\textbf{PCM côté DSP (SOF)} & \textbf{Ch} & \textbf{Rôle / Phase} \\
\midrule
\texttt{SAI\_Capture} & 8 & Dry mics TAC vers A53 (rec/bridge) --- Phase 1 \\
\texttt{SAI\_Playback} & 8 & Monitor mix vers TAC5212 DAC --- Phase 1 \\
\texttt{Master\_Tap} & 2 & Master post-chaîne (NPU + monitor) --- Phase 1 \\
\texttt{USB1\_Send / Ret} & 8 / 8 & Bridge vers UAC2 gadget 1 --- Phase 2 \\
\texttt{USB2\_Send / Ret} & 2 / 2 & Bridge vers UAC2 gadget 2 --- Phase 2 \\
\texttt{FX1..4\_Send / Ret} & 2 / 2 $\times$4 & Bridge vers JACK LV2 --- Phase 3 \\
\bottomrule
\end{tabularx}
\end{center}
\textbf{Total une fois complet~: 11 PCMs}, mais on les ajoute \emph{phase par phase}. Phase 1 = 3 PCMs seulement.

\subsection{Budget latence}
\begin{center}
\begin{tabular}{lrl}
\toprule
\textbf{Étape chemin principal} & \textbf{Latence} & \textbf{Note} \\
\midrule
USB gadget in & 1{,}0~ms & isoc UAC2 @ 48k \\
\texttt{alsaloop} gadget$\rightarrow$DSP (1 période 2~ms) & 2{,}0~ms & single-buffer \\
DSP routing + master chain & 2{,}0~ms & 1 cycle 2ms \\
\texttt{alsaloop} DSP$\rightarrow$gadget & 2{,}0~ms & \\
USB gadget out & 1{,}0~ms & \\
\midrule
\textbf{Round-trip principal} & \textbf{8{,}0~ms} & \textcolor{ca53}{\textbf{$\leq$ 10~ms} \checkmark} \\
\midrule
Chemin mic$\rightarrow$HP (interne DSP pur) & \textbf{$\sim$4~ms} & \\
Chemin avec FX send (DSP$\rightarrow$JACK$\rightarrow$DSP) & $\sim$20--30~ms & non critique \\
\bottomrule
\end{tabular}
\end{center}
\textbf{Contrainte 10~ms tenue} sur le chemin principal. Le mic$\rightarrow$HP traverse uniquement le DSP ($\sim$4~ms), excellent pour monitoring live.

\subsection{Budget CPU}
\begin{center}
\begin{tabular}{lrr}
\toprule
\textbf{Composant} & \textbf{\% HiFi4} & \textbf{\% A53} \\
\midrule
8$\times$ eq\_iir (3 biquads/ch) & 4 & -- \\
Matrice mixin\_mixout & 5--8 & -- \\
multiband\_drc stéréo & 10 & -- \\
volume + drc limiter & 2 & -- \\
Pipelines + IPC SOF & 5 & -- \\
\texttt{alsaloop} $\times$ bridges (USB+FX) & -- & 8--12 \\
JACK + 4 LV2 FX (A53) & -- & 20--30 \\
NPU service + UI & -- & 10--15 \\
\midrule
\textbf{Total estimé} & \textbf{25--30~\%} & \textbf{40--55~\%} \\
\bottomrule
\end{tabular}
\end{center}
Marge confortable côté DSP, charge A53 soutenable ($\sim$1 cœur sur 4 en moyenne).

\subsection{Avantages et inconvénients}
\begin{center}
\begin{tabularx}{\textwidth}{p{0.48\textwidth} X}
\toprule
\textbf{\textcolor{ca53}{Avantages}} & \textbf{\textcolor{chw}{Inconvénients}} \\
\midrule
Contrainte 10~ms \textbf{tenue} & Topology SOF moyenne complexité ($\sim$200 lignes m4) \\
Phase 1 déployable en $\sim$2 semaines (3 PCMs seulement) & Matrice en DSP $=$ ajout de voies $=$ rebuild firmware \\
Effets send dans écosystème LV2 mature & 11 PCMs au total une fois Phase 3 finie \\
CPU DSP $\sim$25~\% $=$ large marge pour exciter/harmonizer futur & \texttt{alsaloop} entre gadget et DSP $=$ 2$\times$2~ms supplémentaires \\
Chaîne master toujours à latence déterministe (DSP) & \\
Codec TAC5212 décharge le DSP du pré-conditionnement & \\
Ajout phasé (mise en route incrémentale sûre) & \\
\bottomrule
\end{tabularx}
\end{center}

\subsection{Verdict Archi C}
\textcolor{ca53}{\textbf{Retenue.}} Seule architecture à tenir la contrainte 10~ms \emph{tout en} gardant un DSP de complexité raisonnable et en réutilisant l'écosystème LV2 pour les effets.

% =======================================================================
\clearpage
\section{Comparaison synthétique}

\subsection{Tableau comparatif}
\begin{center}
\footnotesize
\begin{tabularx}{\textwidth}{lXXX}
\toprule
\textbf{Critère} & \textbf{A --- Max DSP} & \textbf{B --- Min DSP} & \textbf{C --- Équilibrée} \\
\midrule
Matrice routage & DSP \texttt{mixin\_mixout} & JACK (A53) & DSP \texttt{mixin\_mixout} \\
EQ fine 8 entrées & DSP \texttt{eq\_iir} & LV2 (A53) & Codec TAC + DSP léger \\
Chaîne master & DSP (multiband\_drc + vol + drc) & LV2 (A53) & DSP \\
Effets send (Reverb\ldots) & A53 (forcément) & A53 & A53 \\
NPU tap & DSP & A53 JACK capture & DSP \\
PCMs côté SOF & 12--15 & 3 & 3$\rightarrow$11 (phasé) \\
Pipelines SOF & $\sim$15 & 3 & 5$\rightarrow$11 \\
Latence round-trip USB & \textcolor{red!70!black}{\textbf{14~ms}} & $\sim$9--13~ms & \textcolor{ca53}{\textbf{8~ms}} \\
Latence mic$\rightarrow$HP & 4~ms & 13~ms & \textcolor{ca53}{\textbf{4~ms}} \\
CPU DSP & $\sim$35~\% & 3~\% & \textbf{$\sim$25~\%} \\
CPU A53 & $\sim$20~\% & \textcolor{red!70!black}{40--55~\%} & $\sim$40~\% \\
Topology m4 (lignes approx) & \textcolor{red!70!black}{$\sim$500} & \textbf{$\sim$50} & $\sim$200 \\
Ajout d'une voie & rebuild firmware & \textbf{à chaud} & rebuild firmware \\
Débogage quotidien & DSP logs only & \textbf{outils Linux} & mixte \\
Tenue contrainte 10~ms & \textcolor{red!70!black}{non} & \textcolor{red!70!black}{non} & \textcolor{ca53}{\textbf{oui}} \\
\bottomrule
\end{tabularx}
\end{center}

\subsection{Vue d'ensemble du trade-off}

\begin{center}
\begin{tikzpicture}[scale=0.9]
\draw[->, thick] (0,0) -- (12,0) node[right]{\small CPU A53};
\draw[->, thick] (0,0) -- (0,7) node[above]{\small Complexité DSP};

\node[circle, fill=cdsp!30, draw, minimum size=12mm] (A) at (2,6) {A};
\node[circle, fill=ca53!30, draw, minimum size=12mm] (B) at (10,1) {B};
\node[circle, fill=cfx!30, draw, minimum size=12mm] (C) at (6,3.5) {\textbf{C}};

\node[below=6mm of A, font=\scriptsize, align=center]{Max DSP\\{\color{red!70!black}14~ms RTT}};
\node[right=10mm of B, font=\scriptsize, align=center]{Min DSP\\{\color{red!70!black}13~ms mic$\rightarrow$HP}};
\node[below=6mm of C, font=\scriptsize, align=center]{Équilibrée\\{\color{ca53}8~ms RTT}};

% Contour 10ms
\draw[dashed, green!50!black, thick] plot[smooth] coordinates {(0,4.3) (4,4) (7,3) (10,2.5) (12,2.3)};
\node[green!50!black, font=\scriptsize] at (11.5,4){\textit{frontière 10~ms}};

\end{tikzpicture}
\end{center}

\subsection{Décision recommandée}

\textbf{$\rightarrow$ Architecture C (équilibrée)} pour les raisons suivantes~:

\begin{enumerate}[leftmargin=*, itemsep=2pt]
  \item Seule architecture tenant la contrainte \textbf{10~ms round-trip USB}.
  \item Latence mic$\rightarrow$HP la plus basse (4~ms) $=$ monitoring live confortable pour les musiciens.
  \item CPU DSP $\sim$25~\% laisse $\sim$75~\% de marge pour Phase 4 (exciter, harmonizer peut-être) sans rebâtir la topology.
  \item Phase 1 minimale déployable rapidement (3 PCMs $=$ $\sim$100 lignes m4), puis ajout progressif en Phase 2 et 3.
  \item Réutilise la puissance per-channel du TAC5212 (AGC, biquads) qui serait sinon gaspillée.
  \item Garde les effets send dans LV2 (écosystème mature, tunable à chaud).
\end{enumerate}

% =======================================================================
\clearpage
\section{Glossaire}

\begin{description}[leftmargin=2.5cm, style=nextline]
\item[SOF] Sound Open Firmware. Framework firmware open-source qui tourne sur le DSP HiFi4 (Zephyr RTOS comme base) et gère les pipelines audio entre DMA, effets et host.
\item[TAC5212] Codec audio 2 ADC + 4 DAC de TI, avec DSP intégré (biquads, AGC, DRC, HPF) configurable via I\textsuperscript{2}C.
\item[TDM] Time Division Multiplexing. Mode de transport série permettant de véhiculer plusieurs canaux (8 ici) sur une paire BCLK/FSYNC/DATA.
\item[SAI] Synchronous Audio Interface. Contrôleur i.MX8MP gérant le TDM/I\textsuperscript{2}S.
\item[DRC] Dynamic Range Compressor (compresseur).
\item[mixin\_mixout] Composant SOF (API IPC4) permettant un routage N$\times$M avec gain par connexion. Équivalent d'une console de mixage virtuelle.
\item[eq\_iir] Composant SOF d'égaliseur paramétrique à filtres IIR (biquads en cascade).
\item[LV2] Standard de plugins audio Linux, successeur de LADSPA. Plugins tunables à chaud, éditeur graphique via carla/jalv.
\item[JACK] JACK Audio Connection Kit. Serveur audio bas latence pour Linux, graphe de routage dynamique entre applications.
\item[UAC2] USB Audio Class 2. Protocole USB pour transporter de l'audio multi-canaux 24/32-bit jusqu'à 384~kHz.
\item[NPU] Neural Processing Unit (Vivante VIP8000 sur i.MX8MP, 2{,}3~TOPS en int8).
\item[Round-trip latency / RTT] Latence aller-retour entre USB in et USB out côté DAW (ou mic in et HP out).
\item[ASIO] Audio Stream Input/Output. Driver audio bas latence Windows, utilisé par les DAW (Reaper, Ableton, Cubase). Le gadget UAC2 se présente comme un périphérique audio USB que le driver ASIO (ASIO4ALL ou propriétaire) consomme.
\end{description}

\vfill
\begin{center}
\textit{Document généré pour décision d'architecture. Le choix retenu sera reporté dans \texttt{MASTERING\_PLATFORM\_PLAN.md} avant tout code.}
\end{center}

\end{document}
